\chapter{Introducci\'on}

El rendimiento de un motor de aeronave se degrada de forma natural a medida que el equipo envejece y opera en condiciones extremas. El ensuciamiento de las palas (\textit{fouling}), la erosi\'on, la corrosi\'on y el desgaste mec\'anico reducen la eficiencia del motor. Detectar esta degrada-ci\'on a tiempo es importante porque afecta el consumo de combustible, el costo de mantenimiento y las emisiones contaminantes.

La fase de despegue es especialmente cr\'itica. Es el momento en que el motor alcanza su empuje m\'aximo (100 \%), por lo que cualquier p\u00e9rdida de rendimiento se manifiesta con claridad en el flujo de combustible (\textit{Fuel Flow T/O}). Ese par\u00e1metro se traduce directamente en costo, emisiones y necesidad de mantenimiento si no se detecta a tiempo.

El \textit{International Civil Aviation Organization} (ICAO) certifica los motores turbofan y turborreactor mediante el ciclo de aterrizaje y despegue (\textit{Landing and Takeoff}, LTO). Los resultados de esas mediciones se almacenan p\u00fablicamente en el \textit{ICAO Engine Emission Databank} (EEDB). Para cada motor certificado, el EEDB registra variables de dise\u00f1o (tipo de motor, relaci\u00f3n de derivaci\u00f3n, relaci\u00f3n de presi\u00f3n, empuje nominal), las condiciones ambientales de la prueba (presi\u00f3n barom\'etrica, temperatura y humedad) y el flujo de combustible medido en cada r\u00e9gimen~\cite{Kurt2024}.

El flujo de combustible en despegue (\textit{Fuel Flow T/O}) resulta particularmente \u00fatil como indicador de desempe\u00f1o. Al ser el r\u00e9gimen de mayor exigencia t\u00e9rmica y mec\u00e1nica, cualquier desviaci\u00f3n respecto al valor esperado sugiere una p\u00e9rdida de eficiencia. El problema se plantea entonces como una regresi\u00f3n supervisada: estimar el \textit{Fuel Flow T/O} a partir de variables de dise\u00f1o y condiciones ambientales, y comparar el valor real registrado en el EEDB con el valor esperado del modelo.

Este trabajo compara tres m\u00e9todos de aprendizaje autom\u00e1tico: Regresi\'on de Procesos Gaussianos (GPR), M\'aquinas de Vectores de Soporte (SVM) y Perceptr\'on Multicapa (MLP). Los tres modelos se entrenan y validan con los mismos datos del EEDB, para que la comparaci\u00f3n sea justa. El modelo que mejor generaliza se usa para construir intervalos de confianza del 99 \%, que permiten detectar motores con posible deterioro de desempe\u00f1o sin desmontar el motor ni requerir conocimientos expertos.

El objetivo de este documento es desarrollar y validar un modelo basado en aprendizaje autom\u00e1tico que prediga el \textit{Fuel Flow T/O} y detecte deterioros en el rendimiento del motor durante la fase de despegue.
